% Lecture 10 - Page 8 (6-XII-1994).
% Fragment for % Lecture 10 - Page 8 (6-XII-1994).
% Fragment for % Lecture 10 - Page 8 (6-XII-1994).
% Fragment for % Lecture 10 - Page 8 (6-XII-1994).
% Fragment for \input{courses/pages/curs10-pagina8-en}.
% No preamble and no forced page break.
% Continuation of the proof of Lemma 2 from Page 7.
% Faithful transcription note: the manuscript appears to say "a in A_0",
% although A_0 consists of cosets aH. The apparent reading is retained.
\bigskip

\subsection*{Lecture 10 - Page 8}

\noindent $a\in A_0\ \Longleftrightarrow\ haH=aH$, for all $h\in H$
\[
\Longleftrightarrow\ a^{-1}haH=H,\quad(\forall)\,h\in H
\]
\[
\Longleftrightarrow\ a^{-1}Ha=H
\Longleftrightarrow\ H=aHa^{-1}
\Longleftrightarrow\ a\in N_G(H).
\]

\noindent Hence
\[
|A_0|=\left|\frac{N_G(H)}{H}\right|=[N_G(H):H],
\]
\noindent while $|A|=[G:H]$.

\bigskip
\noindent\underline{Corollary}. Under the same hypotheses as Lemma 2,
if $p\mid[G:H]$, then
\[
H\subsetneq N_G(H).
\]

\bigskip
\noindent\fbox{\textbf{Proof of Sylow's First Theorem}}\quad
$p^i\mid |G|$.

\medskip
\noindent Induction on $i$.

\noindent $i=1$: follows from Cauchy's theorem.

\medskip
\noindent Assume $1\leq i<t$ and that there exists $H\leq G$
such that $|H|=p^i$.

\noindent We show that there exists $K\leq G$ with
$|K|=p^{i+1}$ and $H\trianglelefteq K$.

\medskip
\noindent $p\mid[G:H]$, since
\[
|G|=p^t m=[G:H]\,|H|=[G:H]p^i.
\]
\noindent $i<t\ \Longrightarrow\ p^{t-i}\mid[G:H]
\ \Longrightarrow\ p\mid[G:H]$.

\medskip
\noindent By the corollary, $H\subsetneq N_G(H)$.
Consider $N_G(H)/H$.

\medskip
\noindent I claim that
\[
p\mid\left|\frac{N_G(H)}{H}\right|=[N_G(H):H],
\]
\noindent since $p\mid[G:H]$ and
\[
[N_G(H):H]\equiv[G:H]\pmod p.
\]

\noindent By Cauchy's theorem, there exists
$\hat a\in N_G(H)/H$ such that $\operatorname{ord}(\hat a)=p$.

\medskip
\noindent Let $K'\leq N_G(H)/H$, $K'=\langle\hat a\rangle$,
$|K'|=p$.

\[
\varphi:N_G(H)\longrightarrow\frac{N_G(H)}{H},
\qquad \varphi(a)=\hat a.
\]
\[
\varphi\bigl(\varphi^{-1}(K')\bigr)=K'.
\]
% The proof continues on the next manuscript page.
.
% No preamble and no forced page break.
% Continuation of the proof of Lemma 2 from Page 7.
% Faithful transcription note: the manuscript appears to say "a in A_0",
% although A_0 consists of cosets aH. The apparent reading is retained.
\bigskip

\subsection*{Lecture 10 - Page 8}

\noindent $a\in A_0\ \Longleftrightarrow\ haH=aH$, for all $h\in H$
\[
\Longleftrightarrow\ a^{-1}haH=H,\quad(\forall)\,h\in H
\]
\[
\Longleftrightarrow\ a^{-1}Ha=H
\Longleftrightarrow\ H=aHa^{-1}
\Longleftrightarrow\ a\in N_G(H).
\]

\noindent Hence
\[
|A_0|=\left|\frac{N_G(H)}{H}\right|=[N_G(H):H],
\]
\noindent while $|A|=[G:H]$.

\bigskip
\noindent\underline{Corollary}. Under the same hypotheses as Lemma 2,
if $p\mid[G:H]$, then
\[
H\subsetneq N_G(H).
\]

\bigskip
\noindent\fbox{\textbf{Proof of Sylow's First Theorem}}\quad
$p^i\mid |G|$.

\medskip
\noindent Induction on $i$.

\noindent $i=1$: follows from Cauchy's theorem.

\medskip
\noindent Assume $1\leq i<t$ and that there exists $H\leq G$
such that $|H|=p^i$.

\noindent We show that there exists $K\leq G$ with
$|K|=p^{i+1}$ and $H\trianglelefteq K$.

\medskip
\noindent $p\mid[G:H]$, since
\[
|G|=p^t m=[G:H]\,|H|=[G:H]p^i.
\]
\noindent $i<t\ \Longrightarrow\ p^{t-i}\mid[G:H]
\ \Longrightarrow\ p\mid[G:H]$.

\medskip
\noindent By the corollary, $H\subsetneq N_G(H)$.
Consider $N_G(H)/H$.

\medskip
\noindent I claim that
\[
p\mid\left|\frac{N_G(H)}{H}\right|=[N_G(H):H],
\]
\noindent since $p\mid[G:H]$ and
\[
[N_G(H):H]\equiv[G:H]\pmod p.
\]

\noindent By Cauchy's theorem, there exists
$\hat a\in N_G(H)/H$ such that $\operatorname{ord}(\hat a)=p$.

\medskip
\noindent Let $K'\leq N_G(H)/H$, $K'=\langle\hat a\rangle$,
$|K'|=p$.

\[
\varphi:N_G(H)\longrightarrow\frac{N_G(H)}{H},
\qquad \varphi(a)=\hat a.
\]
\[
\varphi\bigl(\varphi^{-1}(K')\bigr)=K'.
\]
% The proof continues on the next manuscript page.
.
% No preamble and no forced page break.
% Continuation of the proof of Lemma 2 from Page 7.
% Faithful transcription note: the manuscript appears to say "a in A_0",
% although A_0 consists of cosets aH. The apparent reading is retained.
\bigskip

\subsection*{Lecture 10 - Page 8}

\noindent $a\in A_0\ \Longleftrightarrow\ haH=aH$, for all $h\in H$
\[
\Longleftrightarrow\ a^{-1}haH=H,\quad(\forall)\,h\in H
\]
\[
\Longleftrightarrow\ a^{-1}Ha=H
\Longleftrightarrow\ H=aHa^{-1}
\Longleftrightarrow\ a\in N_G(H).
\]

\noindent Hence
\[
|A_0|=\left|\frac{N_G(H)}{H}\right|=[N_G(H):H],
\]
\noindent while $|A|=[G:H]$.

\bigskip
\noindent\underline{Corollary}. Under the same hypotheses as Lemma 2,
if $p\mid[G:H]$, then
\[
H\subsetneq N_G(H).
\]

\bigskip
\noindent\fbox{\textbf{Proof of Sylow's First Theorem}}\quad
$p^i\mid |G|$.

\medskip
\noindent Induction on $i$.

\noindent $i=1$: follows from Cauchy's theorem.

\medskip
\noindent Assume $1\leq i<t$ and that there exists $H\leq G$
such that $|H|=p^i$.

\noindent We show that there exists $K\leq G$ with
$|K|=p^{i+1}$ and $H\trianglelefteq K$.

\medskip
\noindent $p\mid[G:H]$, since
\[
|G|=p^t m=[G:H]\,|H|=[G:H]p^i.
\]
\noindent $i<t\ \Longrightarrow\ p^{t-i}\mid[G:H]
\ \Longrightarrow\ p\mid[G:H]$.

\medskip
\noindent By the corollary, $H\subsetneq N_G(H)$.
Consider $N_G(H)/H$.

\medskip
\noindent I claim that
\[
p\mid\left|\frac{N_G(H)}{H}\right|=[N_G(H):H],
\]
\noindent since $p\mid[G:H]$ and
\[
[N_G(H):H]\equiv[G:H]\pmod p.
\]

\noindent By Cauchy's theorem, there exists
$\hat a\in N_G(H)/H$ such that $\operatorname{ord}(\hat a)=p$.

\medskip
\noindent Let $K'\leq N_G(H)/H$, $K'=\langle\hat a\rangle$,
$|K'|=p$.

\[
\varphi:N_G(H)\longrightarrow\frac{N_G(H)}{H},
\qquad \varphi(a)=\hat a.
\]
\[
\varphi\bigl(\varphi^{-1}(K')\bigr)=K'.
\]
% The proof continues on the next manuscript page.
.
% No preamble and no forced page break.
% Continuation of the proof of Lemma 2 from Page 7.
% Faithful transcription note: the manuscript appears to say "a in A_0",
% although A_0 consists of cosets aH. The apparent reading is retained.
\bigskip

\subsection*{Lecture 10 - Page 8}

\noindent $a\in A_0\ \Longleftrightarrow\ haH=aH$, for all $h\in H$
\[
\Longleftrightarrow\ a^{-1}haH=H,\quad(\forall)\,h\in H
\]
\[
\Longleftrightarrow\ a^{-1}Ha=H
\Longleftrightarrow\ H=aHa^{-1}
\Longleftrightarrow\ a\in N_G(H).
\]

\noindent Hence
\[
|A_0|=\left|\frac{N_G(H)}{H}\right|=[N_G(H):H],
\]
\noindent while $|A|=[G:H]$.

\bigskip
\noindent\underline{Corollary}. Under the same hypotheses as Lemma 2,
if $p\mid[G:H]$, then
\[
H\subsetneq N_G(H).
\]

\bigskip
\noindent\fbox{\textbf{Proof of Sylow's First Theorem}}\quad
$p^i\mid |G|$.

\medskip
\noindent Induction on $i$.

\noindent $i=1$: follows from Cauchy's theorem.

\medskip
\noindent Assume $1\leq i<t$ and that there exists $H\leq G$
such that $|H|=p^i$.

\noindent We show that there exists $K\leq G$ with
$|K|=p^{i+1}$ and $H\trianglelefteq K$.

\medskip
\noindent $p\mid[G:H]$, since
\[
|G|=p^t m=[G:H]\,|H|=[G:H]p^i.
\]
\noindent $i<t\ \Longrightarrow\ p^{t-i}\mid[G:H]
\ \Longrightarrow\ p\mid[G:H]$.

\medskip
\noindent By the corollary, $H\subsetneq N_G(H)$.
Consider $N_G(H)/H$.

\medskip
\noindent I claim that
\[
p\mid\left|\frac{N_G(H)}{H}\right|=[N_G(H):H],
\]
\noindent since $p\mid[G:H]$ and
\[
[N_G(H):H]\equiv[G:H]\pmod p.
\]

\noindent By Cauchy's theorem, there exists
$\hat a\in N_G(H)/H$ such that $\operatorname{ord}(\hat a)=p$.

\medskip
\noindent Let $K'\leq N_G(H)/H$, $K'=\langle\hat a\rangle$,
$|K'|=p$.

\[
\varphi:N_G(H)\longrightarrow\frac{N_G(H)}{H},
\qquad \varphi(a)=\hat a.
\]
\[
\varphi\bigl(\varphi^{-1}(K')\bigr)=K'.
\]
% The proof continues on the next manuscript page.
